% =====================================================
% 题型：长方体截面周长填空题 (q_solid_cuboid_section.tex)
% =====================================================
% =====================================================
% 难度：★★★ (难题)
% =====================================================
\newcommand{\qSolidCuboidSectionStem}{如图，在长方体 \(ABCD-A_1B_1C_1D_1\) 中，\(AD = AA_1 = \dfrac{1}{2} AB = 2\)，\(\vec{A_1P} = \dfrac{3}{4}\vec{AB} + \dfrac{1}{4}\vec{A_1D_1}\)，则过点 \(C, D_1, P\) 的平面截长方体的截面周长为 \underline{\hspace{2.5cm}}。}

\newcommand{\qSolidCuboidSectionFigure}[1][1.3]{%
  \begin{tikzpicture}[scale=#1, >=stealth, line join=round, line cap=round]
    % 1. 严格使用向量基底生成坐标，确保透视绝对平行
    \coordinate (A)  at (0,0);
    \coordinate (B)  at (3.2,0);
    \coordinate (D)  at (0.9,0.55);
    \coordinate (A1) at (0,1.8);
    
    \coordinate (C)  at ($(B)+(D)-(A)$);
    \coordinate (B1) at ($(B)+(A1)-(A)$);
    \coordinate (C1) at ($(C)+(A1)-(A)$);
    \coordinate (D1) at ($(D)+(A1)-(A)$);

    % 2. 优化 P 点计算逻辑：
    % 因 \vec{AB} = \vec{A_1B_1}，原式即为 \vec{A_1P} = 3/4\vec{A_1B_1} + 1/4\vec{A_1D_1}
    % 这意味着 P 位于线段 D1B1 上，可直接调用 TikZ 比例计算
    \coordinate (P) at ($(D1)!0.75!(B1)$);

    % 3. 绘制内部不可见线段（被遮挡的边）
    \draw[dashed, semithick] (A) -- (D) -- (C);
    \draw[dashed, semithick] (D) -- (D1);

    % 4. 教师端解析涂色与高亮辅助线优化（严格遵守遮挡关系）
    \ifteacher
      \fill[blue!15, opacity=0.8] (C) -- (D1) -- (B1) -- cycle;
      \draw[dashed, thick, blue!80] (C) -- (D1);
      \draw[thick, blue!80] (D1) -- (B1) -- (C);
      \draw[very thick, red!70] (D1) -- (B1);
      \node[blue!80, inner sep=1.5pt] at ($(C)!0.5!(D1)!0.18!(B1)$)
        {\scriptsize \(\triangle CB_1D_1\)};
    \fi

    % 5. 绘制外部可见实体轮廓
    \draw[thick] (A) -- (B) -- (C) -- (C1) -- (D1) -- (A1) -- cycle;
    \draw[thick] (B) -- (B1) -- (C1);
    \draw[thick] (A1) -- (B1);

    % 6. 已知条件 A1P 与点 P
    % P 位于可见的上底面，从 A1 连向 P 的线段代表已知向量 A1P
    \draw[densely dashed, thick] (A1) -- (P);
    \fill (P) circle (1.5pt) node[above left, inner sep=1.5pt] {\small $P$};

    % 7. 顶点标签精调位置（避免与线条交叉重叠）
    \node[below left]  at (A)  {\small $A$};
    \node[below]       at (B)  {\small $B$};
    \node[right]       at (C)  {\small $C$};
    \node[left]        at (D)  {\small $D$};
    \node[left]        at (A1) {\small $A_1$};
    \node[above right] at (B1) {\small $B_1$};
    \node[right]       at (C1) {\small $C_1$};
    \node[above left]  at (D1) {\small $D_1$};
  \end{tikzpicture}%
}

\newcommand{\qSolidCuboidSectionAnswer}{\(4\sqrt{5} + 2\sqrt{2}\)}

\newcommand{\qSolidCuboidSectionSolution}{%
  以 \(A\) 为原点，\(AB, AD, AA_1\) 所在射线分别为 \(x, y, z\) 轴的正方向，建立空间直角坐标系 \(A-xyz\). \\
  由已知 \(AD = AA_1 = 2\), \(AB = 4\)，可得相关点坐标为：
  \[
    A(0,0,0), \quad B(4,0,0), \quad D(0,2,0), \quad C(4,2,0),
  \]
  \[
    A_1(0,0,2), \quad B_1(4,0,2), \quad D_1(0,2,2), \quad C_1(4,2,2).
  \]
  根据 \(\vec{A_1P} = \dfrac{3}{4}\vec{AB} + \dfrac{1}{4}\vec{A_1D_1}\)，且 \(\vec{AB} = (4,0,0)\), \(\vec{A_1D_1} = (0,2,0)\)，得：
  \[
    \vec{A_1P} = \frac{3}{4}(4,0,0) + \frac{1}{4}(0,2,0) = (3, 0.5, 0).
  \]
  所以点 \(P\) 的坐标为 \(P(3, 0.5, 2)\). \\
  因为 \(\vec{D_1P} = (3, -1.5, 0)\)，\(\vec{D_1B_1} = (4, -2, 0)\)，
  易知 \(\vec{D_1P} = \dfrac{3}{4}\vec{D_1B_1}\)， \\
  所以点 \(P\) 位于上底面对角线 \(D_1B_1\) 上，
  从而过点 \(C, D_1, P\) 的平面即为过点 \(C, D_1, B_1\) 的平面，该平面截长方体所得截面为 \(\triangle C B_1 D_1\). \\
  利用空间两点间距离公式，计算三角形各边长为：
  \[
    B_1D_1 = \sqrt{(4-0)^2 + (0-2)^2 + (2-2)^2} = \sqrt{20} = 2\sqrt{5},
  \]
  \[
    D_1C = \sqrt{(0-4)^2 + (2-2)^2 + (2-0)^2} = \sqrt{20} = 2\sqrt{5},
  \]
  \[
    CB_1 = \sqrt{(4-4)^2 + (2-0)^2 + (0-2)^2} = \sqrt{8} = 2\sqrt{2}.
  \]
  所以，该平面截长方体的截面周长为：
  \[
    L = B_1D_1 + D_1C + CB_1 = 2\sqrt{5} + 2\sqrt{5} + 2\sqrt{2} = 4\sqrt{5} + 2\sqrt{2}.
  \]
  故填 \(4\sqrt{5} + 2\sqrt{2}\).%
}

\newcommand{\qSolidCuboidSectionDiagnosis}{%
  若得到类似 \(8+\sqrt{13}\) 的结果，通常说明截面识别有误：应先判断 \(P\) 位于上底面对角线 \(B_1D_1\) 上，从而得到截面为 \(\triangle CB_1D_1\)，再计算三条截面边长。%
}

\newcommand{\qSolidCuboidSectionTeachingNotes}{%
  必讲。课堂主线为“坐标定点 \(P\) -> 判断 \(P,D_1,B_1\) 共线 -> 确定截面 \(\triangle CB_1D_1\) -> 三维距离公式求周长”。教师端图要突出 \(B_1D_1, CB_1, CD_1\)。%
}
